\chapter{Transformation Operators}\label{cha:operators}

This chapter describes the four truth-value-preserving transformation operators implemented in the QBF Transformer tool. For each operator we give a formal definition, a proof of truth-value preservation, and a description of the expected effect on solver behavior.

All operators follow the same contract: given an input formula $\varphi$, the operator produces a \emph{mutant} $\varphi'$ such that $\varphi$ is true if and only if $\varphi'$ is true. The original formula object is never modified; the operator always works on a deep copy.

\section{Tautology Insertion}\label{sec:tautology}

\subsection{Definition}

Let $\varphi$ be a QBF formula and let $x$ be an existentially quantified variable in $\varphi$. The \emph{tautology insertion} operator produces the mutant
\[
  \varphi' = \varphi \land (x \lor \lnot x).
\]
The clause $(x \lor \lnot x)$ is called a \emph{tautology clause} because it evaluates to true under every possible assignment to $x$.

\subsection{Proof of Truth-Value Preservation}

\begin{proof}
Let $\varphi$ be any QBF formula and $x$ any Boolean variable.
The clause $(x \lor \lnot x)$ is a tautology, i.e.\ it is true for every assignment to $x$.
Therefore:
\[
  \varphi \land (x \lor \lnot x) \;\equiv\; \varphi \land \top \;\equiv\; \varphi.
\]
The truth value of $\varphi'$ is identical to that of $\varphi$ under every quantifier assignment.
\end{proof}

\subsection{Expected Effect on Solver}

The tautology clause adds one syntactic clause to the formula without adding any logical constraint. 
Both DepQBF and CAQE must parse and process the additional clause during preprocessing. 
Because the clause contains a complementary literal pair $(x, \lnot x)$, a solver that performs tautology detection during clause simplification will immediately recognize and discard it. 
The expected runtime overhead is therefore small (ratio close to $1.0$ and slightly below), and the coverage delta is expected to be low, as no fundamentally new solver code paths are activated.

\begin{listing}[H]
\caption{Tautology insertion example.}
\label{lst:tautology}
\begin{minted}[frame=single,fontsize=\small,breaklines]{text}
% Original (3 variables, 2 clauses)
p cnf 3 2
a 1 0
e 2 3 0
1 -2 0
-2 3 0

% Mutant: one tautology clause added for variable 2
p cnf 3 3
a 1 0
e 2 3 0
1 -2 0
-2 3 0
2 -2 0
\end{minted}
\end{listing}

\section{Pure Literal Generator}\label{sec:pureliteral}

\subsection{Definition}

A variable $x$ is a \emph{pure literal} in a formula $\varphi$ if it appears in $\varphi$ with only one polarity: either exclusively as $x$ (positive pure literal) or exclusively as $\lnot x$ (negative pure literal), but never both.

The \emph{pure literal generator} operator identifies an existentially quantified variable $x$ that is already a pure literal in $\varphi$ and applies \emph{Pure Literal Elimination} (PLE) to produce the mutant $\varphi'$:

\begin{itemize}
  \item If $x$ is a positive pure literal: set $x = \top$. Remove all clauses containing $x$ (they are satisfied). Remove the variable $x$ from the quantifier prefix.
  \item If $x$ is a negative pure literal: set $x = \bot$. Remove all clauses containing $\lnot x$ (they are satisfied). Remove $x$ from the prefix.
\end{itemize}

\subsection{Proof of Truth-Value Preservation}

\begin{proof}
Let $x$ be an existentially quantified positive pure literal in $\varphi$, i.e.\ $\lnot x$ does not appear in any clause of $\varphi$.
Because $\lnot x$ is absent, assigning $x = \top$ satisfies every clause containing $x$ and leaves all other clauses unchanged. No clause is falsified by this assignment. The existential player can always choose $x = \top$, so the truth value of $\varphi$ is not affected by the elimination.
The argument for a negative pure literal is symmetric.
\end{proof}

\textbf{Remark.} The operator only applies PLE to variables that are \emph{already} pure literals in the original formula. Artificially making a variable pure by removing occurrences of one polarity would not be truth-value-preserving in general.

\subsection{Expected Effect on Solver}

Pure literal elimination removes clauses from the formula, making the 
mutant strictly smaller than the original. The solver is expected to 
solve the mutant faster (runtime ratio $< 1.0$). Because the mutant 
is smaller, the solver traverses fewer code paths overall, resulting 
in a negative coverage delta relative to the original run.

\section{Universal Reduction Trigger}\label{sec:ur}

\subsection{Definition}
\emph{Universal Reduction} (UR) is a simplification rule specific to QBF~\cite{KleineBuening1995}: a universal literal $u$ may be removed from a clause $C$ if no existential variable in $C$ lies at a deeper quantifier level than $u$'s block. Formally, let $\mathrm{level}(v)$ denote the index of the quantifier block containing variable $v$ in the prefix, with level $0$ corresponding to the outermost block. Then $u$ is removable from $C$ if:
\[
  \forall e \in C,\; e \text{ existential} \;\Rightarrow\; \mathrm{level}(e) < \mathrm{level}(u).
\]

The \emph{universal reduction trigger} operator inserts the clause
\[
  C_{\mathrm{UR}} = [u,\; e,\; \lnot e]
\]
where $u$ is a universal variable and $e$ is an existential variable with $\mathrm{level}(u) < \mathrm{level}(e)$, i.e.\ $u$ lies at an outer level than $e$. This level configuration is chosen so that the solver, when inspecting the clause, has to consult the same quantifier-level and clause-simplification machinery that universal reduction relies on, even though the clause is ultimately discarded as a tautology.

\subsection{Proof of Truth-Value Preservation}

\begin{proof}
The clause $[u, e, \lnot e]$ is a tautology for every assignment to $e$:
\[
  [u, e, \lnot e] =
  \begin{cases}
    [u, \top, \bot] = \top & \text{if } e = \top,\\
    [u, \bot, \top] = \top & \text{if } e = \bot.
  \end{cases}
\]
The clause evaluates to true regardless of the value of $u$ or $e$. Adding a tautological clause to any formula preserves its truth value:
\[
  \varphi \;\equiv\; \varphi \land \top \;\equiv\; \varphi \land [u, e, \lnot e].
\]
\end{proof}

\textbf{Remark on a prior incorrect design.} An earlier version of this operator inserted the clause $[u, e]$ instead of $[u, e, \lnot e]$. This is \emph{not} truth-value-preserving: if the universal player assigns $u = \bot$, the clause becomes a unit clause $[e]$, which may force $e = \top$ and create a conflict with existing clauses, potentially turning a satisfiable formula into an unsatisfiable one. The tautological form $[u, e, \lnot e]$ avoids this issue entirely.

\subsection{Expected Effect on Solver}

When the solver encounters $C_{\mathrm{UR}}$, it must inspect the quantifier levels of $u$ and $e$ during its standard simplification phase: the clause contains a universal literal whose level is outer to that of the existential pair $e, \lnot e$. Even though the clause is ultimately discarded as a tautology (because of the pair $[e, \lnot e]$), the solver traverses the same level-checking and clause-simplification code that is also used during universal reduction. The runtime overhead is therefore expected to be small, while the coverage delta should be slightly positive --- reflecting that some preprocessing code paths are reached that the original formula does not exercise.

\section{Blocked Clause Insertion}\label{sec:blocked}

\subsection{Definition}

Let $C$ be a clause and $l \in C$ a literal. $C$ is \emph{blocked} with respect to $l$ if, for every clause $D$ in the formula that contains $\lnot l$, the resolvent $C \otimes_l D = (C \setminus \{l\}) \cup (D \setminus \{\lnot l\})$ is a tautology~\cite{BiereEtAl2011}.

The \emph{blocked clause insertion} operator constructs such a clause $C$ and adds it to $\varphi$. The simplest construction arises when no clause in $\varphi$ contains $\lnot l$: in this case the condition holds vacuously for every clause, and the single-literal clause $[l]$ is trivially blocked. More generally, the operator selects an existential literal $l$ and builds $C$ by appending literals that ensure every resolvent with a $\lnot l$-clause is a tautology.
\subsection{Proof of Truth-Value Preservation}

\begin{proof}
This follows directly from the Blocked Clause Addition (BCA) theorem~\cite{Jarvisalo2010,BiereEtAl2011}: if $C$ is blocked with respect to some literal $l \in C$, then
\[
  \varphi \;\equiv\; \varphi \land C.
\]
Intuitively, a blocked clause can always be satisfied by setting $l = \top$ without falsifying any other clause, because every clause $D$ containing $\lnot l$ already contains a tautological resolvent pair with $C$. The truth value of $\varphi$ is therefore unchanged. 
\end{proof}

\subsection{Expected Effect on Solver}

Both DepQBF and CAQE implement blocked clause detection as part of their preprocessing phase. Inserting a blocked clause forces the solver to execute the BCE/BCA detection code paths, which increases coverage for the corresponding solver modules. The runtime effect depends on the size of the formula and the cost of BCE detection; in general, the ratio is expected to be close to $1.0$ with a moderate positive coverage delta.

\section{Summary}\label{sec:operator-summary}

Table~\ref{tab:operators} summarises the four operators, their structural change to the formula, the basis of truth-value preservation, and the primary solver mechanism each is designed to activate.

\begin{table}[h]
\centering
\caption{Overview of the four transformation operators.}
\label{tab:operators}
\resizebox{\textwidth}{!}{%
\begin{tabular}{llll}
\hline
\textbf{Operator} & \textbf{Structural change} & \textbf{Preservation basis} & \textbf{Target mechanism} \\
\hline
Tautology Insertion         & Adds clause $[x, \lnot x]$         & Tautology ($\equiv \top$) & Tautology detection       \\
Pure Literal Generator      & Removes clauses, shrinks prefix     & PLE soundness             & Pure Literal Elimination  \\
Universal Reduction Trigger & Adds clause $[u, e, \lnot e]$      & Tautology ($\equiv \top$) & Universal Reduction       \\
Blocked Clause Insertion    & Adds blocked clause $C$             & BCA theorem               & Blocked Clause Elimination\\
\hline
\end{tabular}}
\end{table}

%%% Local Variables:
%%% tex-command: "pdflatex"
%%% TeX-master: "../main"
%%% TeX-command-extra-options: "-shell-escape"
%%% End: